\section{Resultados experimentales y análisis de rendimiento}

\subsection{Comparación de FID}

El ponente mostró una comparación de la distancia de Fréchet Inception (FID) entre distintos métodos de muestreo en conjuntos de datos como CIFAR-10:

\begin{itemize}
    \item \textbf{DDPM} (1000 pasos): FID aproximadamente 3$\sim$4 (línea base)
    \item \textbf{DDIM} (100 pasos): FID con degradación notable, aproximadamente 5$\sim$8
    \item \textbf{DDIM} (10 pasos): FID empeora considerablemente, aproximadamente 20+
    \item \textbf{DPM-Solver-2} (15$\sim$20 pasos): FID cercano al nivel de DDPM en 1000 pasos
    \item \textbf{DPM-Solver-3} (10$\sim$15 pasos): FID alcanza e incluso supera al de DDPM en 1000 pasos
\end{itemize}

\begin{importantbox}{Resultados clave}
Utilizando DPM-Solver-3, solo se requieren 15 pasos para alcanzar la calidad de muestreo de DDPM en 1000 pasos. Esto implica una aceleración de aproximadamente \textbf{70 veces}, y \textbf{sin necesidad de volver a entrenar el modelo} —logrado completamente mediante un solucionador ODE más eficiente durante la fase de inferencia—. Este es también uno de los puntos destacados que el ponente enfatiza reiteradamente: distintas perspectivas matemáticas (SDE $\to$ ODE $\to$ estructura semilineal $\to$ integradores exponenciales) pueden traducirse directamente en mejoras del rendimiento práctico.
\end{importantbox}

\subsection{Comportamiento de convergencia según el orden}

\begin{knowledgebox}{Relación entre orden y número de pasos}
\begin{itemize}
    \item Cuando hay suficientes pasos ($>50$), las diferencias entre los distintos DPM-Solver son menores, ya que el error inherente al método de Euler ya es muy pequeño.
    \item Con pocos pasos ($10\sim 20$), la ventaja de los métodos de orden superior es significativa; el segundo orden mejora un orden de magnitud respecto al primero.
    \item Con muy pocos pasos ($<10$), el tercer orden aún ofrece mejoras sobre el segundo, aunque con rendimientos decrecientes.
    \item Por encima del tercer orden, el error de aproximación intrínseco de la red ($\boldsymbol{\epsilon}_\theta \neq \boldsymbol{\epsilon}^*$) se convierte en un cuello de botella.
\end{itemize}
\end{knowledgebox}

\subsection{Relación con otros métodos de aceleración}

DPM-Solver pertenece a los métodos de aceleración de muestreo \textbf{Training-free} (sin entrenamiento adicional). Complementarios a este enfoque están:

\begin{table}[H]
\centering
\begin{tabular}{lcc}
\toprule
\textbf{Método} & \textbf{¿Requiere entrenamiento?} & \textbf{Enfoque central} \\
\midrule
DPM-Solver/DDIM & No & Mejor discretización de ODE \\
Destilación progresiva (Progressive Distillation) & Sí & Un modelo estudiante imita a un maestro multietapa \\
Modelos de consistencia (Consistency Models) & Sí & Aprendizaje directo del mapeo ruido $\to$ datos \\
Flow Matching + Euler & No & Trayectorias más rectas; el método de Euler resulta más efectivo \\
\bottomrule
\end{tabular}
\caption{Comparación de distintas estrategias de aceleración de muestreo}
\end{table}

\begin{knowledgebox}{Complementariedad entre DPM-Solver y Flow Matching}
Flow Matching entrena para que las trayectorias ODE sean más rectas (cercanas a lineales), permitiendo que el método de Euler funcione muy bien. Por su parte, DPM-Solver busca soluciones numéricas más eficientes sobre trayectorias ODE dadas (posiblemente curvas). Ambos enfoques son ortogonales y complementarios: si la trayectoria ODE ya es suficientemente recta, la ventaja de DPM-Solver disminuye; sin embargo, para las trayectorias curvas típicas de los modelos difusivos estándar, el valor de DPM-Solver es muy significativo.
\end{knowledgebox}

\subsection{Resumen de la sección}

Los experimentos demuestran que DPM-Solver comprime el número de pasos de muestreo desde 1000 hasta 10$\sim$20 sin volver a entrenar el modelo, logrando aceleraciones de varias decenas de veces. Su ventaja central radica en aprovechar la estructura semilineal de las ODEs de los modelos difusivos, algo que los solucionadores genéricos de ODE no pueden hacer.

